\section{总结与延伸}

这堂课从 AGI 的抽象形态出发，最终落到一套可执行的系统地图。Agent 的核心不是“模型会调用工具”，而是让目标、状态、动作、权限、评测、学习、记忆、通信和恢复形成闭环。任何一层缺失，都会把能力演示变成不可维护的自动化。

更进一步说，本讲把“模型智能”与“系统可信”明确分开：前者决定候选动作的上限，后者决定哪些动作能进入真实世界。一个更强模型可以减少部分推理错误，却不会自动提供审计、授权、回滚、版本和责任人。学习 Agent 系统时，应始终把模型改进放回完整控制面中评价。

\subsection{一张表串起整套 Agent stack}

\begin{center}
\begin{tabularx}{\textwidth}{>{\bfseries}p{0.16\textwidth} X X X}
层次 & 核心问题 & 代表机制 & 失败信号 \\
\hline
Goal & 用户真正要什么 & 约束解析、确认 & 目标漂移、隐含约束丢失 \\
Policy & 下一步做什么 & planning、MCTS、routing & 贪心、循环、越权 \\
Environment & 动作在哪里执行 & browser/API/tools & 页面变化、不可逆动作 \\
Evaluation & 怎样知道做对 & REAL、slice、trace checks & 总分掩盖尾部 \\
Learning & 怎样从轨迹改进 & critique、DPO、RL & 奖励投机、污染数据 \\
Memory & 怎样跨步骤与会话保存状态 & scratchpad、retrieval、profile & 过期、泄露、错误持久化 \\
Coordination & 多个 Agent 怎样合作 & hierarchy、sync、MCP/A2A & 冲突、重复、等待 \\
Deployment & 怎样承担现实责任 & budget、audit、override & 失控、无追踪、无法恢复 \\
\end{tabularx}
\end{center}

\begin{importantbox}{本讲最重要的压缩}
AGI 的产品形态尚未确定，但可靠 Agent 的工程约束已经很清楚：\textbf{可复现地评、在轨迹上学、按权限行动、带来源记忆、通过协议协作、在失败时停止并交还控制。}
\end{importantbox}

\subsection{开始一个 Agent 项目前的十二问}

\begin{multicols}{2}
\small
\begin{enumerate}
  \item 用户目标和完成条件能否机器判定？
  \item 哪些动作可逆，哪些需要确认？
  \item Agent 通过 API 还是直接控制界面？
  \item benchmark 是否覆盖真实网站、任务长度和风险 slice？
  \item 是否记录完整 observation--action--result trace？
  \item 失败 taxonomy 能否定位到模型、工具、memory 或 policy？
  \item 搜索和 self-critique 是否有成本与停止条件？
  \item preference data 是否同时包含成功和失败分支？
  \item 长期记忆是否有来源、过期、权限和用户编辑？
  \item 多 agent 是否真的可并行，还是只增加通信？
  \item MCP/A2A 之外，身份、授权和版本怎样管理？
  \item 上线后谁监控、谁接管、怎样回滚？
\end{enumerate}
\normalsize
\end{multicols}

\subsection{拓展阅读}

\begin{itemize}
  \item Putta et al., \emph{Agent Q: Advanced Reasoning and Learning for Autonomous AI Agents}：MCTS、self-critique 与偏好学习。\href{https://arxiv.org/abs/2408.07199}{arXiv:2408.07199}
  \item AGI, Inc., \emph{REAL: Benchmarking Autonomous Agents on Deterministic Simulations of Real Websites}：确定性网页环境与 leaderboard。\href{https://github.com/agi-inc/real}{REAL repository}
  \item Lilian Weng, \emph{LLM Powered Autonomous Agents}：memory、planning、tools 与 action 的架构综述。\href{https://lilianweng.github.io/posts/2023-06-23-agent/}{article}
  \item Model Context Protocol 官方文档：工具、资源与 prompt 的开放连接协议。\href{https://modelcontextprotocol.io/}{official docs}
\end{itemize}

\paragraph{证据边界。}
本讲记录的是 2025 年 4 月的课堂内容。REAL、AgentQ、MultiOn、MCP 和当时的模型结果都应按课堂时间点理解；讲义保留它们的机制与工程启示，不把 demo、产品状态或 benchmark 数字改写成 2026 年的当前产品保证。

\end{document}
